% =====================================================================
% Software Development  --  TARGET: ~700 words
% Marking weight: 15 marks. THIS IS THE HIGHEST-LEVERAGE SECTION FOR
% BREAKING THE 60% CAP. The 1st+ rubric explicitly names ISRs,
% FreeRTOS, state machines, structured protocols.
% =====================================================================
\section{Software Development}
\label{sec:software}

\guide{
This is where most of your advanced-feature claims should land. The
1st+ rubric wants ``highly efficient, modular and well-documented
code with real-time performance optimisations, robust error handling.''
Six sub-sections, in this order:
\begin{enumerate}
  \item Main control loops (one paragraph each for Controller and
        Arm).
  \item Wireless protocol (custom 7-byte frame --- this is your
        biggest single advanced claim; include the listing).
  \item Closed-loop joint control (PID, ISR-driven).
  \item Operational modes \& FSM (autonomous-mode state diagram).
  \item Vision co-processor integration.
  \item Documentation \& testing.
\end{enumerate}
Use the listings package to embed code excerpts (already loaded in
main.tex). Don't copy-paste big blocks --- 5--15 lines max per
listing, with comments.
}

\subsection{Main control loops}
\placeholder{[\ldots one paragraph for the Controller's super-loop:
joystick read $\rightarrow$ dead-band $\rightarrow$ scale to angle
$\rightarrow$ pack into frame $\rightarrow$ TX at 50\,Hz. One
paragraph for the Arm's super-loop: UART servicing, 100\,Hz PID per
joint, vision UART, FSM update, 500\,ms link-loss watchdog.\ldots]}

\subsection{Wireless protocol}
\label{sec:protocol}
\guide{The custom packet format is a strong advanced claim. Show:
\begin{itemize}
  \item The frame layout as a table or inline diagram.
  \item Why each field is what it is (start byte for re-sync; XOR
        for cheap single-bit error detection on AVR; 1-byte angles
        because servo resolution is the bottleneck).
  \item Bit-rate calculation: $7\,\text{B}\times 10\,\text{bits/B}\times 50\,\text{Hz} = 3.5\,\text{kbit/s}$
        well inside HC-05's 9600 baud.
  \item Failure modes: bad checksum $\rightarrow$ drop frame; no
        valid frame for 500\,ms $\rightarrow$ park in safe pose.
\end{itemize}
Then embed the encoder/decoder listing.}

\placeholder{[\ldots one or two paragraphs explaining the frame format
and rationale\ldots]}

\begin{lstlisting}[caption={Frame encoder/decoder (excerpt from
\texttt{lib/ArmProtocol/ArmProtocol.h}).},label=lst:proto]
inline void encode(const Frame& f, uint8_t out[FRAME_SIZE]) {
    out[0] = START_BYTE;
    out[1] = f.base;     out[2] = f.shoulder;
    out[3] = f.elbow;    out[4] = f.claw;
    out[5] = f.flags;    out[6] = checksum(f);
}

bool Decoder::feed(uint8_t b, Frame& out) {
    buf_[idx_++] = b;
    if (idx_ == 1 && buf_[0] != START_BYTE) { idx_ = 0; return false; }
    if (idx_ < FRAME_SIZE) return false;
    idx_ = 0;
    Frame f { buf_[1], buf_[2], buf_[3], buf_[4], buf_[5] };
    return checksum(f) == buf_[6] ? (out = f, true) : false;
}
\end{lstlisting}

\subsection{Closed-loop joint control}
\label{sec:closedloop}
\guide{Show: the standard discrete PID equation; the 100\,Hz timer
ISR; anti-windup clamp on the integrator; what was tuned and how.
A step-response plot ($K_p$ alone, then PI, then PID) is a strong
1st-class artefact --- include even a rough one.}

\placeholder{[\ldots one paragraph on the PID and tuning
methodology\ldots]}

% TODO FIGURE: step-response plot from serial logger
\begin{figure}[t]
  \centering
  \fbox{\parbox{0.9\columnwidth}{\centering\vspace{4pt}
  \emph{[Step-response plot: commanded vs measured angle on one
  joint, showing $K_p$, PI and full-PID traces.]}\vspace{4pt}}}
  \caption{Closed-loop step response (placeholder).}
  \label{fig:stepresp}
\end{figure}

\subsection{Operational modes and finite state machine}
\label{sec:fsm}
\guide{Draw the FSM in TikZ or import as PDF. The states required by
the brief: IDLE $\rightarrow$ SCAN $\rightarrow$ DETECT $\rightarrow$
CLASSIFY $\rightarrow$ PICK $\rightarrow$ PLACE $\rightarrow$
COUNT\_CHECK, with a SAFE state covering checksum-fail, link-loss
and e-stop. Mention the $\geq$5-item brief target. The FSM is named
in the 1st+ rubric --- don't skip the diagram.}

\placeholder{[\ldots one paragraph describing the FSM, listing each
state and its transition conditions, and how mode switching from
manual is gated\ldots]}

% TODO FIGURE: autonomous-mode state diagram (TikZ or PDF)
\begin{figure}[t]
  \centering
  \fbox{\parbox{0.9\columnwidth}{\centering\vspace{4pt}
  \emph{[State diagram: IDLE / SCAN / DETECT / CLASSIFY / PICK /
  PLACE / COUNT\_CHECK / SAFE.]}\vspace{4pt}}}
  \caption{Autonomous-mode finite state machine (placeholder).}
  \label{fig:fsm}
\end{figure}

\subsection{Vision co-processor integration}
\guide{Brief paragraph saying the vision co-processor publishes a
one-line ASCII record \texttt{CAM,<state>,<colour>,<pct>$\backslash$n}
to the Mega's Serial1, and that the autonomous-mode FSM consumes this
record in the CLASSIFY state. Forward-reference
Section~\ref{sec:vision}.}

\placeholder{[\ldots one paragraph on the camera UART protocol and
how the FSM consumes it\ldots]}

\subsection{Documentation and testing}
\guide{Mention: code organised as a PlatformIO mono-repo with three
build environments; protocol library is host-testable; in-tree unit
tests cover encode/decode round-trip + checksum rejection +
re-synchronisation. List the number of tests passing.}

\placeholder{[\ldots short paragraph on code organisation, comments,
and host-side unit tests\ldots]}

\corehit{Brief Objectives 1, 2, 5 and~6: digital command packets
captured at the controller, transmitted wirelessly, received and
acted on at the arm; manual and autonomous modes both implemented.}

\advhit{[\ldots e.g.\ custom 7-byte framed UART protocol with
start-byte re-synchronisation, XOR checksum, link-loss watchdog;
ISR-driven 100\,Hz PID with anti-windup; explicit FSM with SAFE
fault state; host-side unit tests for the protocol\ldots]}
